% Ad Atticum 13.39 --- headnote
% ad-atticum-13-39, 5 August 45 BC, Tusculum

\textit{Cicero to Atticus, written at the Tusculan villa around 5
August 45 BC --- Perseus dateline \textit{Scr. in Tusculano m.
Non. Sext. a. 709 (45)}. Two very short sections, the third
letter in three days about the nephew. Cicero opens with a
single hot exclamation --- ``What incredible falsity!'' --- on
learning that young Quintus tells his father one story and his
mother the opposite. Then a turn: the fellow is already going
soft, conceding that his father has every right to be angry.}

\textit{The second section concedes Atticus's tactical advice from
13.38: the \textit{crooked path} it shall be. Cicero will come to
Rome, reluctantly, dragging himself away from the
\textit{Academica}. He notes Brutus's return from the latest of
his journeys ``not from the place I would have wished'' --- the
glancing reference is probably to a visit to Caesar's circle ---
and asks Atticus for books he needs, including a treatise by
Phaedrus the Epicurean, \textit{On the Gods}. A daggered crux at
the very end (a missing second title in the manuscript tradition)
is preserved in place. The register is hurried: short clauses,
mid-sentence shifts, the inside-baseball ``Brutus too,'' you
say --- the back-and-forth of a daily correspondence.}
